\documentclass[10pt,a4paper]{amsart}
\usepackage[margin=19mm]{geometry}
\usepackage{amsmath,amssymb,mathtools}
\usepackage[scr=rsfs,cal=euler]{mathalfa}
\usepackage{xfrac}
\usepackage[unicode,hidelinks,bookmarks=false]{hyperref}
\newcommand{\ie}{i.e.\ }
\newcommand{\cf}{cf.\ }
\newcommand{\resp}{respectively\ }
\newcommand{\Subsection}{\subsection}
\providecommand{\nobreakdash}{\nobreak\hbox{-}}
\setlength{\emergencystretch}{2em}
\multlinegap0pt
\usepackage{physics}
\usepackage{mathtools}
\usepackage{mathrsfs}
\usepackage{relsize}
\usepackage[dvipsnames]{xcolor}
\usepackage{enumitem}
\usepackage{cleveref}
\numberwithin{equation}{section}   %%%%%%%%%%%%%%
\allowdisplaybreaks
%%%%%%%%    Theorems etc.   %%%%%%%%%%%%%%%%%%%%%
\theoremstyle{plain}   %%%%%%%%%%%%%%%%%%%%%%%%%%
\newtheorem{thm}{Theorem}   %%%%%%%%%%%%%%%%%%%%%
\renewcommand{\thethm}{\arabic{thm}} 
\newtheorem{leftemma}{Lemma}
\newtheorem{manualtheoreminner}{Theorem}    %%%%%
\newenvironment{manualtheorem}[1]{% %%%%%%%%%%%%%
  \renewcommand\themanualtheoreminner{#1}%  %%%%%
  \manualtheoreminner       %%%%%%%%%%%%%%%%%%%%%
}{\endmanualtheoreminner}   %%%%%%%%%%%%%%%%%%%%%
\crefname{thm}{theorem}{theorems}
\theoremstyle{plain}   %%%%%%%%%%%%%%%%%%%%%%%%%%
\newtheorem{assumption}{Assumption}     %%%%%%%%%
\crefname{assumption}{assumption}{assumptions}
\newtheorem{manualainner}{Assumption}   %%%%%%%%%

\newenvironment{manuala}[1]{%       %%%%%%%%%%%%%
  \renewcommand\themanualainner{#1}%        %%%%%
  \manualainner         %%%%%%%%%%%%%%%%%%%%%%%%%
}{\endmanualainner}     %%%%%%%%%%%%%%%%%%%%%%%%%
%%%%%%%%%%%%%%%%%%%%%%%%%%%%%%%%%%%%%%%%%%%%%%%%%

%%%%%%%%%%%%%%%%%%%%%%%%%%%%%%%%%%%%%%%%%%%%%%%%%
\newtheorem{lemma}{Lemma}[section]   %%%%%%%%%%%%
\crefname{lemma}{lemma}{lemmas}
\newtheorem{theorem}{Theorem}[section]  %%%%%%%%%
\newtheorem{bigcor}[thm]{Corollary}
\crefname{bigcor}{corollary}{corollaries}
\Crefname{bigcor}{Corollary}{Corollaries}
\newtheorem{bigprop}[thm]{Proposition}
\newtheorem{definition}{Definition}[section]    %
\newtheorem{prop}[lemma]{Proposition}   %%%%%%%%%
\newtheorem{cor}[lemma]{Corollary}   %%%%%%%%%%%%
\newtheorem{remark}[lemma]{Remark}   %%%%%%%%%%%%
\newtheorem{dfn}{Definition}    %%%%%%%%%%%%%%%%%   
%%%%%%%%    Colord QED Symbole              %%%%%
\renewcommand{\qedsymbol}{\rule{0.7em}{0.7em}} %%
%%%%%%%%%%%%%%%%%%%%%%%%%%%%%%%%%%%%%%%%%%%%%%%%%
\newtheoremstyle{examplestyle}%
  {3pt}{3pt}% Space above/below
  {\normalfont}% Body font
  {}% Indent
  {\bfseries}% Head font
  {.}% Punctuation after head
  {0.5em}% Space after head
  {}% Head spec
\theoremstyle{examplestyle}
\newtheorem{example}{Example}
\newcommand{\exend}{\unskip\nobreak\hfill\ensuremath{\scalebox{1.3}{$\diamond$}}\par}
%%%%%%%%%%%%%%%%%%%%%%%%%%%%%%%%%%%%%%%%%%%%%%%%%%%%%%%%%%%%%%%%%%%%%%%%%%%%
\newtheorem{condition}{Condition}
\crefname{condition}{condition}{conditions}
\Crefname{condition}{Condition}{Conditions}
\providecommand{\theHcondition}{\arabic{condition}}
\newcommand{\condtag}[1]{\renewcommand{\thecondition}{(#1)}}
\newcommand{\condreset}{\renewcommand{\thecondition}{\arabic{condition}}}
%%%%%%%%%%%%%%%%%%%%%%%%%%%%%%%%%%%%%%%%%%%%%%%%%%%%%%%%%%%%%%%%%%%%%%%%%%%%
%%%%%%%%    Math notation syntax                    %%%%%%%%%%%%%%%%%%%%%%%%
\ExplSyntaxOn       %%%%%%%%%%%%%%%%%%%%%%%%                             %%%
\NewDocumentCommand{\definealphabet}{mmmm}{%                             %%%
  \int_step_inline:nnn { `#3 } { `#4 }{%                                 %%%
    \cs_new_protected:cpx { #1 \char_generate:nn { ##1 }{ 11 } }{%       %%%
      \exp_not:N #2 { \char_generate:nn { ##1 } { 11 } }}}}              %%%           
\ExplSyntaxOff      %%%%%%%%%%%%%%%%%%%%%%%%                             %%%
\definealphabet{mb}{\mathbb}{A}{Z}      %%%%%%%%                         %%%  
\definealphabet{mc}{\mathcal}{A}{Z}     %%%%%%%%                         %%%
\definealphabet{scr}{\mathscr}{A}{Z}    %%%%%%%%                         %%%
\definealphabet{mf}{\mathfrak}{A}{Z}    %%%%%%%%                         %%%
%%%%%%%%%%%%%%%%%%%%%%%%%%%%%%%%%%%%%%%%%%%%%%%%%%%%%%%%%%%%%%%%%%%%%%%%%%%%

%%%%%%%%%%%%%%%%%%%%%%%%%%%%%%%%%%%%%%%%%%%%%%%%%%%%%%%%%%%%%%%%%%%%%%%%%%%%
%%%%%%%%    Standalone Math Operators          %%%%%%%%%%%%%%%%%%%%%%%%%%%%%
\DeclareMathOperator{\states}{\mbS_d}   %%%%%%%%%   states amthbb{S}_d  %%%%
\DeclareMathOperator{\tens}{\mbT_{d,D}}
\DeclareMathOperator{\mapspace}{\mcL(\mbM_d)}   %%%%%%%%%%%%%%%%%%%%%%%%%%%%
\DeclareMathOperator{\vari}{\text{var}}         %%%%%%%%%%%%%%%%%%%%%%%%%%%%
\DeclarePairedDelimiterX{\corr}[2]{\mathrm{corr}(}{)}{#1,#2}     %%%%%%%%%%%
\DeclarePairedDelimiterX{\cov}[2]{\mathrm{cov}(}{)}{#1,#2}          %%%%%%%%
\DeclarePairedDelimiterX{\dis}[2]{\mathrm{d}(}{)}{#1, #2} %%%%%%%%%%%%%%%%%%
\DeclareMathOperator{\probspace}{(\Omega, \mcF, \pr)}   %%%%%%%%%%%%%%%%%%%%
\DeclareMathOperator{\ergsys}{(\Omega, \mcF, \pr, \{\theta_g\}_{g\in\mbR})}
\DeclareMathOperator{\dummyborel}{\mathscr{B}}                  %%%%%%%%%%%%
\newcommand\borel[1]{\dummyborel\left(#1\right)}    %%%%%%%%%%%%%%%%%%%%%%%%
\DeclareMathOperator{\pr}{\mbP}           %%%%%%%%%%%%%%%%%%%%%%%%%%%
\newcommand{\adj}{^\dagger}         %%%%%%%%%%%%%%%%%%%%%%%%%%%%%%%%%%%%%%%%
\DeclareMathOperator{\outcomes}{\mcA^\mbN}      %%%%%%%%%%%%%%%%%%%%%%%%%%%%
\DeclareMathOperator{\matrices}{\mbM_d}     %%%%%%%%%%%%%%%%%%%%%%%%%%%%%%%%
\DeclarePairedDelimiterX{\inner}[2]{\langle}{\rangle}{#1| #2}   %%%%%%%%%%%%
\DeclareMathOperator*{\esssup}{ess\,sup}            %%%%%%%%%%%%%%%%%%%%%%%%
\DeclareMathOperator*{\essinf}{ess\,inf}        %%%%%%%%%%%%%%%%%%%%%%%%%%%%
\DeclareMathOperator{\ten}{\mbC^d\otimes(\mbC^D)^{\otimes 2}}       %%%%%%%%
\newcommand{\law}[1]{\pr-\text{Law}[#1]}            %%%%%%%%%%%%%%%%%%%%%%%%    
\newcommand{\event}{E_*}            %%%%%%%%%%%%%%%%%%%%%%%%%%%%%%%%%%%%%%%%
\newcommand{\Law}{\mathrm{Law}}     %%%%%%%%%%%%%%%%%%%%%%%%%%%%%%%%%%%%%%%%
\newcommand{\s}{\rho_{\mathsf{ss}}}     %%%%%%%%%%%%%%%%%%%%%%%%%%%%%%%%%%%%
\newcommand{\Dynerg}{{\hyperlink{Dyn-Erg}{(Dyn-Erg) }}}     %%%%%%%%%%%%%%%%
\newcommand{\dee}{\mathrm{d}}           %%%%%%%%%%%%%%%%%%%%%%%%%%%%%%%%%%%%
\newcommand{\Diag}{\operatorname{Diag}}     %%%%%%%%%%%%%%%%%%%%%%%%%%%%%%%%
\newcommand{\Deph}{\mathsf{Deph}}           %%%%%%%%%%%%%%%%%%%%%%%%%%%%%%%%
\newcommand{\AD}{\mathsf{AD}}       %%%%%%%%%%%%%%%%%%%%%%%%%%%%%%%%%%%%%%%%
%%%%%%%%%%%%%%%%%%%%%%%%%%%%%%%%%%%%%%%%%%%%%%%%%%%%%%%%%%%%%%%%%%%%%%%%%%%%

%%%%%%%%%%%%%%%%%%%%%%%%%%%%%%%%%%%%%%%%%%%%%%%%%%%%%%%%%%%%%%%%%%%%%%%%%%%%
%%%%%%%%%%%%%%%%%%%%%%%%%%%%%%%%%%%%%%%%%%%%%%%%%%%%%%%%%%%%%%%%%%%%%%%%%%%%
\newcommand{\esp}{{\hyperlink{esp}{ESP}}}
\newcommand{\rhomix}{{\hyperlink{rho_mix}{$\rho$-mixing}}}
%%%%%%%%%%%%%%%%%%%%%%%%%%%%%%%%%%%%%%%%%%%%%%%%%%%%%%%%%%%%%%%%%%%%%%%%%%%%
%%%%%%%%%%%%%%%%%%%%%%%%%%%%%%%%%%%%%%%%%%%%%%%%%%%%%%%%%%%%%%%%%%%%%%%%%%%%

%%%%%%%%%%%%%%%%%%%%%%%%%%%%%%%%%%%%%%%%%%%%%%%%%%%%%%%%%%%%%%%%%%%%%%%%%%%%
%%%%%%%%    Defining trace  %%%%%%%%%%%%%%%%%%%%%%%%%%%%%%%%%%%%%%%%%%%%%%%%
\let\oldtr\tr %%%%%%%%%%%%%%%%%%%%%%%%%%%%%%%%%%%   tr{} scales %%%%%%%%%%%%
\renewcommand{\tr}[1]{\operatorname{Tr}\left(#1\right)} %   tr[] sclaes %%%%
\newcommand{\otr}[1]{\operatorname{Tr}\left[#1\right]}
\DeclareMathOperator{\str}{tr} %%%%%%%%%%%%%%%%%%   tr          %%%%%%%%%%%%
\newcommand{\ptr}[2]{{\operatorname{tr}_{#1}\left[#2\right]}}   %%%%%%%%%%%%
%%%%%%%%%%%%%%%%%%%%%%%%%%%%%%%%%%%%%%%%%%%%%%%%%%%%%%%%%%%%%%%%%%%%%%%%%%%%
\newcommand{\oln}{\overline}                        %%%%%%%%%%%%%%%%%%%%%%%%
\DeclareMathOperator{\bQ}{\oln{\mbQ}}               %%%%%%%%%%%%%%%%%%%%%%%%
\newcommand{\normalize}[1]{\frac{#1}{\tr{#1}}}      %%%%%%%%%%%%%%%%%%%%%%%%
%%%%%%%%%%%%%%%%%%%%%%%%%%%%%%%%%%%%%%%%%%%%%%%%%%%%%%%%%%%%%%%%%%%%%%%%%%%%

\renewcommand{\vec}[1]{\text{vec}(#1)}

%%%%%%%%%%%%%%%%%%%%%%%%%%%%%%%%%%%%%%%%%%%%%%%%%%%%%%%%%%%%%%%%%%%%%%%%%%%%
%%%%%%%%    Defining Rank-1 Sup-op      %%%%%%%%%%%%%%%%%%%%%%%%%%%%%%%%%%%%
\newcommand{\opp}[2]{\Xi_{(#1, #2)}}        %%%%%%%%%%%%%%%%%%%%%%%%%%%%%%%%
\newcommand{\wopp}[2]{\widetilde\Xi_{(#1, #2)}}         %%%%%%%%%%%%%%%%%%%%
%%%%%%%%%%%%%%%%%%%%%%%%%%%%%%%%%%%%%%%%%%%%%%%%%%%%%%%%%%%%%%%%%%%%%%%%%%%%

%%%%%%%%%%%%%%%%%%%%%%%%%%%%%%%%%%%%%%%%%%%%%%%%%%%%%%%%%%%%%%%%%%%%%%%%%%%%
\newcommand{\avg}[1]{\mbE\left[#1\right]}           %%%%%%%%%%%%%%%%%%%%%%%%
\DeclarePairedDelimiterX{\cnum}[1]{\mathrm{c}(}{)}{#1}  %%%%%%%%%%%%%%%%%%%%
\newcommand{\floor}[1]{\left\lfloor #1 \right\rfloor}   %%%%%%%%%%%%%%%%%%%%
\newcommand{\ceil}[1]{\left\lceil #1 \right\rceil}      %%%%%%%%%%%%%%%%%%%%
\DeclarePairedDelimiterX{\diam}[1]{\text{diam}(}{)}{#1} %%%%%%%%%%%%%%%%%%%%
\newcommand{\proj}{{\mathlarger{\boldsymbol{\cdot}}}} %%%%%%%%%%%%%%%%%%%%%%
\let\oldker\ker         %%%%%%%%%%%%%%%%%%%%%%%%%%%%%%%%%%%%%%%%%%%%%%%%%%%%
\renewcommand{\ker}[1]{\oldker\left(#1\right)}  %%%%%%%%%%%%%%%%%%%%%%%%%%%%
\let\oldtilde\tilde     %%%%%%%%%%%%%%%%%%%%%%%%%%%%%%%%%%%%%%%%%%%%%%%%%%%%
\renewcommand{\tilde}{\widetilde}       %%%%%%%%%%%%%%%%%%%%%%%%%%%%%%%%%%%%
%%%%%%%%%%%%%%%%%%%%%%%%%%%%%%%%%%%%%%%%%%%%%%%%%%%%%%%%%%%%%%%%%%%%%%%%%%%%

%%%%%%%%%%%%%%%%%%%%%%%%%%%%%%%%%%%%%%%%%%%%%%%%%%%%%%%%%%%%%%%%%%%%%%%%%%%%
%%%%%%%%%%%%%%%%%%%%%%%%%%%%%%%%%%%%%%%%%%%%%%%%%%%%%%%%%%%%%%%%%%%%%%%%%%%%

% Notation for the additional functional and contraction statements.
\DeclareMathOperator{\Cov}{Cov}
\DeclareMathOperator{\Var}{Var}
\DeclareMathOperator{\id}{id}
\newcommand{\TV}{d_{\mathrm{TV}}}
\newcommand{\ind}{\mathbf1}
\newcommand{\unorm}[1]{\left\lVert#1\right\rVert}
\title[Central limit theorems for outcome records]{Central Limit Theorems for Outcome Records in Disordered Quantum Trajectories: the second-moment comparison for arbitrary initial states (Proposition 5.4)}
\author[Lubashan Pathirana]{Lubashan Pathirana}
\date{Source: arXiv:2603.28893v2 (CC BY 4.0)}
\begin{document}
\maketitle
\vspace{1em}
\noindent\textit{Source and licence.} Central Limit Theorems for Outcome Records in Disordered Quantum Trajectories. Version arXiv:2603.28893v2 (CC BY 4.0), distributed under the Creative Commons Attribution 4.0 International licence, http://creativecommons.org/licenses/by/4.0/, as recorded in the arXiv OAI licence metadata for this exact version and shown on the abstract page https://arxiv.org/abs/2603.28893v2. This item is an editorial re-presentation of the paper's Proposition 5.4 together with the paper's own proof of it. A full rights notice is delivered at licenses/arxiv-2603.28893-CC-BY-4.0.txt.
\noindent\textit{Editorial scope.} Counted unit: the paper's second-moment comparison for an arbitrary measurable initial state, printed as Proposition 5.4 (source0.tex lines 4139--4150), delivered together with the paper's complete original proof of it (source0.tex lines 4151--4547) as the single primary proof body. No mathematics is written, repaired or re-derived: statement and proof are the paper's own text line for line, in the paper's own notation ($\overline\mbQ^{\mathrm{out}}_\vartheta$, $\overline\mbQ_{\s}$, $\beta_n(\vartheta)$, $\Phi^{(n)}_\omega$, $\unorm{\cdot}_1$, $\mcA^\ell$, $\states$). Required context retained uncounted: the paper's Assumption 1 (source0.tex lines 568--571), which the counted statement invokes by cross-reference, and the heading of the paper's Section 5. Every cross-reference inside the retained spans resolves inside the delivered document: the only reference command of statement or proof is the citation-free cross-reference to Assumption 1, and no citation command occurs in any retained span. Printed numbers are the paper's own: the wrapper restores the paper's section counter, its shared section-relative lemma/proposition counter and its section-relative equation counter before the retained material, so the retained section prints as Section 5, the retained assumption as Assumption 1, the counted statement as Proposition 5.4, and the two numbered displays of the proof as (5.5) and (5.6). Omitted from this package and not redistributed: the paper's preamble and front matter as such (source0.tex lines 1--235 and 236--261, of which only the title and the author name are reproduced in the wrapper); the examples, the preliminaries, the stationary functional CLT proofs and the rest of the universal CLT section, whose statements and proofs lie outside the retained ranges; and the paper's own bibliography machinery with its bib data file. The delivered document therefore carries no bibliography, and no reference command of the delivered text was rewritten or adapted: there was nothing to adapt. Scope deviation, disclosed: the counted unit is a proposition of the paper's universal-CLT section rather than its headline theorem, because the headline theorem's proof is cited through a chain that leaves the citation-free region of the paper, and a faithful bibliography cannot be rebuilt from a biblatex-only source without redistributing or inventing reference data. The counted proposition is the second-moment half of the paper's universality argument and its proof is the paper's own 397-line comparison of the coupled posterior states under two different initial laws; the whole argument is delivered. Printed slips kept literal, among them the paper's own wording and its own choice of symbols. Layout and class: the paper's own class is the standard article class with the a4paper option, which is not a publisher class, and it is not shipped with this package; the delivered wrapper is the lane's pinned amsart editorial wrapper at 10pt on A4, to which this package adds the paper's own theorem-like declarations with their shared counter, its own notation macros and its own operator declarations, all copied verbatim from the paper's preamble. No publisher class file, no publisher style file and no upstream script is redistributed. Assets: no retained span of this package contains a figure environment, a drawing picture or a graphics-inclusion command, no image or artwork file exists in or is redistributed with this package, and the manifest and the delivery evidence both carry an empty asset-dependency list. A full rights notice is delivered at licenses/arxiv-2603.28893-CC-BY-4.0.txt.
\setcounter{section}{4}
\section{Universal CLT}

\setcounter{lemma}{3}
\setcounter{equation}{4}
    \begin{assumption}
    \label{assumption1}
        The base system \((\Omega,\mcF,\pr,\theta)\) is invertible, ergodic, and \(\pr\)-preserving.
    \end{assumption}

\begin{prop}
\label{univ:prop:moments}
    Suppose \Cref{assumption1} holds and $\s$ is a dynamically stationary state. 
    Assume the uniform forgetting bound $\beta_n(\eta)\le r_n\to0$ for all measurable initial states $\eta:\Omega\to\states$, and $\sum_{n\ge0}\beta_n(\vartheta)<\infty$ for the given measurable initial state $\vartheta:\Omega\to\states$. 
    For every bounded scalar $g:\mcA^\ell\to\mbR$ with $\mbE_{\overline\mbQ^{\mathrm{out}}_{\s}} g(A_{1:\ell})=0$, put $X_i=g(A_{i:i+\ell-1})$ and $S_n=\sum_{i=1}^nX_i$. Then
    \[
     \mbE_{ \overline\mbQ^{\mathrm{out}}_\vartheta} S_n=O(1),
     \qquad \text{ and }\qquad
     \frac1n\left(\mbE_{ \overline\mbQ^{\mathrm{out}}_\vartheta} S_n^2-\mbE_{ \overline\mbQ^{\mathrm{out}}_{\s}} S_n^2\right)\longrightarrow0.
    \]
    Consequently, any stationary second-moment limit is also the second-moment and covariance limit under $\overline\mbQ_\vartheta$ after division by $n$.
\end{prop}

\begin{proof}
    Let $g:\mcA^\ell\to\mbR$ be a bounded scalar function with  $\mbE_{\s}g(A_{1:\ell})=0$, and define
    \[
        B:=\max_{w\in\mcA^\ell}|g(w)|.
    \]
    For the random initial state $\vartheta$, put $b_i:=\beta_{i-1}(\vartheta)$, so $\sum_{i\ge1}b_i<\infty$.
    Throughout this proof, we retain the shorthand $\mbE_\vartheta :=\mbE_{\overline\mbQ^{\mathrm{out}}_\vartheta}$ and $\mbE_{\s} :=\mbE_{\overline\mbQ^{\mathrm{out}}_{\s}}$ for expectations under the corresponding annealed outcome laws.
    Set
    \[
        \Delta_{i-1}(\omega)
        :=
        \Phi^{(i-1)}_\omega(\vartheta(\omega))
        -
        \s(\theta^{i-1}\omega).
    \]
    For $j\ge1$, let
    \[
        H_j(\omega)
        :=
        \sum_{w\in\mcA^\ell}
            g(w)
            \bigl(\mcT^{[j:j+\ell-1]}_{w;\omega}\bigr)^*
            (\mbI_d),
    \]
    where the star denotes the Hilbert--Schmidt adjoint.
    This is the Hermitian operator representing $g$ on the block starting at $j$.
    The block effects are positive and sum to $\mbI_d$, so
    \[
        -B\mbI_d\le H_j(\omega)\le B\mbI_d,
        \qquad
        \unorm{H_j(\omega)}_\infty\le B.
    \]
    Summing over the first $i-1$ outcomes gives
    \[
        \mbE_\vartheta X_i-\mbE_{\s}X_i
        =
        \mbE_{\pr}
            \str\bigl(H_i\Delta_{i-1}\bigr).
    \]
    Consequently,
    \[
        |\mbE_\vartheta X_i-\mbE_{\s}X_i|
        \le
        B\mbE_{\pr}\unorm{\Delta_{i-1}}_1
        =
        Bb_i.
    \]
    Stationarity of the annealed outcome law under $\s$ gives $\mbE_{\s}X_i=0$ for every $i$. 
    Hence
    \[
        |\mbE_\vartheta S_n|
        \le
        B\sum_{i=1}^n b_i
        \le
        B\sum_{i\ge1}b_i<\infty.
    \]
    The upper bound is finite and independent of $n$.
    Thus $\mbE_\vartheta S_n=O(1)$, proving the first assertion.

    For $0\le j-i<\ell$, the product $X_iX_j$ is a single finite-block observable bounded in absolute value by $B^2$.
    Applying the same argument to this observable shows that
    \[
        \left|
            \mbE_\vartheta[X_iX_j]
            -
            \mbE_{\s}[X_iX_j]
        \right|
        \le B^2b_i.
    \]

    Since $X_iX_j=X_jX_i$, the diagonal terms appear once in the second-moment difference, while the upper-triangular terms appear twice. For each fixed $i$, the overlapping upper-triangular terms have
    \[
        i+1\le j\le\min(n,i+\ell-1).
    \]
    Using the preceding bound, we get that the quantity 
    \[
        \frac1n\sum_{i=1}^n
            \left|
                \mbE_\vartheta[X_i^2]
                -
                \mbE_{\s}[X_i^2]
            \right|
        +
        \frac2n\sum_{i=1}^n
        \sum_{j=i+1}^{\min(n,i+\ell-1)}
            \left|
                \mbE_\vartheta[X_iX_j]
                -
                \mbE_{\s}[X_iX_j]
            \right|
    \]
    is bounded above by 
    \[
        \frac{B^2}{n}\sum_{i=1}^n b_i
        +
        \frac{2(\ell-1)B^2}{n}\sum_{i=1}^n b_i
        =
        \frac{(2\ell-1)B^2}{n}\sum_{i=1}^n b_i
        \le
        \frac{2\ell B^2}{n}\sum_{i\ge1}b_i.
    \]
    But $\frac{2\ell B^2}{n}\sum_{i\ge1}b_i\to 0$ as $n\to\infty$. 

    For the nonoverlapping terms, define at time $i+\ell-1$ the Hermitian matrix
    \[
        D_i(\omega)
        :=
        \sum_{w\in\mcA^\ell}
            g(w)
            \mcT^{[i:i+\ell-1]}_{w;\omega}
                (\Delta_{i-1}(\omega)),
    \]
    and take $u_i:=\str D_i$. 
    For any Hermitian $D$, the Jordan decomposition and trace preservation of the sum of the block operations give
    \[
        \begin{aligned}
            \sum_w\unorm{\mcT_w(D)}_1
            &\le
            \sum_w
                \left(
                    \str\mcT_w(D_+)
                    +
                    \str\mcT_w(D_-)
                \right)\\
            &=
            \sum_w\str\mcT_w(|D|)
            =
            \unorm D_1,
        \end{aligned}
    \]
    where $\mcT_w$ abbreviates  $\mcT^{[i:i+\ell-1]}_{w;\omega}$.
    Hence
    \[
        \unorm{D_i}_1
        \le B\unorm{\Delta_{i-1}}_1
        \le2B,
        \qquad
        \mbE_{\pr}\unorm{D_i}_1\le Bb_i.
    \]

    Fix $j\ge i+\ell$ and put $h:=j-i-\ell$.
    Define
    \[
        v_j(\omega)
        :=
        \str\bigl(
            H_j(\omega)\s(\theta^{j-1}\omega)
        \bigr)
    \]
    and
    \[
        R_{ij}
        :=
        \mbE_{\pr}\str\left[
            H_j
            \left(
                \Phi^{(h)}_{\theta^{i+\ell-1}\omega}(D_i)
                -
                u_i\s(\theta^{j-1}\omega)
            \right)
        \right].
    \]
    Born's rule now gives the exact difference
    \begin{align}
    \label{univ:eq:pairdifference}
        \mbE_\vartheta[X_iX_j]
        -
        \mbE_{\s}[X_iX_j]
        &=
        \mbE_{\pr}\str\bigl(
            H_j
            \Phi^{(h)}_{\theta^{i+\ell-1}\omega}(D_i)
        \bigr)
        \nonumber\\
        &=
        \mbE_{\pr}[u_i v_j]+R_{ij}.
    \end{align}
    The number of discarded channels is exactly $h$: their times are $i+\ell,\ldots,j-1$.
    The product is the identity when $h=0$.

    Since $|u_i|\le\unorm{D_i}_1$ and positive trace-preserving maps contract the trace norm of
    Hermitian matrices,
    \[
        \mbE_{\pr}
        \unorm{
            \Phi^{(h)}_{\theta^{i+\ell-1}\omega}(D_i)
            -
            u_i\s(\theta^{j-1}\omega)
        }_1
        \qquad\le
        2\mbE_{\pr}\unorm{D_i}_1
        \le2Bb_i.
    \]

    A second bound follows by writing $D_i=D_i^+-D_i^-$ and setting
    \[
        c_i^\pm:=\str D_i^\pm,
        \qquad
        \eta_i^\pm:=
        \begin{cases}
            D_i^\pm/c_i^\pm, & c_i^\pm>0,\\
            \rho_*, & c_i^\pm=0,
        \end{cases}
    \]
    where $\rho_*\in\states$ is a fixed state.
    Positive and negative parts are continuous functions of a Hermitian matrix, so $\eta_i^\pm$ are measurable states.
    We have
    \[
        \Phi^{(h)}_{\theta^{i+\ell-1}\omega}(D_i)
        -
        u_i\s(\theta^{j-1}\omega)
        =
        c_i^+
        \left(
            \Phi^{(h)}_{\theta^{i+\ell-1}\omega}(\eta_i^+)
            -
            \s(\theta^{j-1}\omega)
        \right)
        -
        c_i^-
        \left(
            \Phi^{(h)}_{\theta^{i+\ell-1}\omega}(\eta_i^-)
            -
            \s(\theta^{j-1}\omega)
        \right).
    \]
    
    Each random mass satisfies $c_i^\pm\le\unorm{D_i}_1\le2B$ pointwise.
    Change variables $\xi=\theta^{i+\ell-1}\omega$ and apply uniform forgetting to the measurable states
    \[
        \xi\longmapsto
        \eta_i^\pm\bigl(\theta^{-(i+\ell-1)}\xi\bigr).
    \]
    Since $\theta$ preserves $\pr$ and $j-1=i+\ell-1+h$, this gives
    \[
        \mbE_{\pr}
        \unorm{
            \Phi^{(h)}_{\theta^{i+\ell-1}\omega}(\eta_i^\pm)
            -
            \s(\theta^{j-1}\omega)
        }_1
        \le r_h.
    \]
    Bounding the random masses before taking expectations therefore yields
    \[
        \mbE_{\pr}
        \unorm{
            \Phi^{(h)}_{\theta^{i+\ell-1}\omega}(D_i)
            -
            u_i\s(\theta^{j-1}\omega)
        }_1
        \le4Br_h.
    \]
    Together with $\unorm{H_j}_\infty\le B$, the two bounds imply
    \begin{equation}
    \label{univ:eq:remainder}
        |R_{ij}|
        \le
        \min\bigl(2B^2b_i,\,4B^2r_h\bigr).
    \end{equation}
    These estimates remain valid at $h=0$ with $r_0=2$.

    For $n,i\ge1$, put
    \[
        a_{n,i}:=\frac1n\sum_{j=i+\ell}^n|R_{ij}|,
    \]
    with empty sums understood to be zero.
    By \eqref{univ:eq:remainder},
    \[
        0\le a_{n,i}
        \le
        \frac{4B^2}{n}\sum_{h=0}^{n-1}r_h
        \longrightarrow0,
    \]
    because $r_h\to0$.
    Moreover,
    \[
        a_{n,i}\le2B^2b_i,
        \qquad
        \sum_{i\ge1}b_i<\infty.
    \]
    Dominated convergence for the series over $i$ gives
    \[
        \frac1n
        \sum_{i=1}^n\sum_{j=i+\ell}^n|R_{ij}|
        =
        \sum_{i\ge1}a_{n,i}
        \longrightarrow0.
    \]

    It remains to handle $\mbE_{\pr}[u_i v_j]$, where $u_i$ may depend on the entire environment.
    Define the bounded measurable function
    \[
        v(\omega)
        :=
        \str\bigl(H_1(\omega)\s(\omega)\bigr).
    \]
    The definition of the block operators gives $H_j=H_1\circ\theta^{j-1}$, and hence
    \[
        v_j=v\circ\theta^{j-1},
        \qquad
        |v_j|\le B,
        \qquad
        \mbE_{\pr}v=\mbE_{\s}X_1=0.
    \]
    The mean ergodic theorem for the unitary operator $U\psi:=\psi\circ\theta$ on $L^2(\pr)$ gives convergence of the Cesaro averages to the orthogonal projection onto the invariant functions.
    Ergodicity makes this projection the constant $\mbE_{\pr}v=0$. 
    Thus
    \[
        \frac1n\sum_{j=1}^n v_j
        \longrightarrow0
        \qquad\text{in }L^2(\pr),
    \]
    and hence also in $L^1(\pr)$.

    For each fixed $i$, removing the first $i+\ell-1$ terms preserves this convergence.
    Since $|u_i|\le2B$, it follows that
    \[
            \left|
                \frac1n
                \sum_{j=i+\ell}^n
                    \mbE_{\pr}[u_i v_j]
            \right|
            =
            \left|
                \mbE_{\pr}\left[
                    u_i\frac1n\sum_{j=i+\ell}^n v_j
                \right]
            \right|
            \le
            2B\mbE_{\pr}
                \left|
                    \frac1n\sum_{j=i+\ell}^n v_j
                \right|
            \to0.
    \]
    Also, for every $n,i\ge1$,
    \[
        \left|
            \frac1n
            \sum_{j=i+\ell}^n\mbE_{\pr}[u_i v_j]
        \right|
        \le
        B\mbE_{\pr}|u_i|
        \le B^2b_i.
    \]
    Another dominated convergence argument over $i$ yields
    \[
        \frac1n
        \sum_{i=1}^n\sum_{j=i+\ell}^n
            \mbE_{\pr}[u_i v_j]
        \longrightarrow0.
    \]

    To combine these estimates, write
    \[
        \delta_{ij}
        :=
        \mbE_\vartheta[X_iX_j]
        -
        \mbE_{\s}[X_iX_j].
    \]
    Since the observables are scalar,
    \[
        \frac1n
        \bigl(
            \mbE_\vartheta S_n^2-\mbE_{\s}S_n^2
        \bigr)
        =
        \frac1n\sum_{i=1}^n\delta_{ii}
        +
        \frac2n\sum_{i=1}^n\sum_{j=i+1}^n\delta_{ij}.
    \]
    The diagonal and overlapping contributions tend to zero by the earlier estimate. 
    The nonoverlapping contribution tends to zero by \eqref{univ:eq:pairdifference} and the two dominated convergence arguments above.
    This proves the second-moment comparison.
    
    Finally,
    \[
            \frac1n
            \left(
                \operatorname{Var}_{\overline\mbQ_\vartheta}(S_n)
                -
                \mbE_{\s}S_n^2
            \right)
            =
            \frac1n
            \left(
                \mbE_\vartheta S_n^2-\mbE_{\s}S_n^2
            \right)
            -
            \frac{(\mbE_\vartheta S_n)^2}{n}
            \longrightarrow0,
    \]
    since $\mbE_\vartheta S_n=O(1)$.
    This proves the variance assertion whenever the stationary second-moment limit exists.
\end{proof}

\end{document}
